%======================================================================
\section{Method}
\label{sec:method}
%======================================================================

This section defines the model, states the structural properties that
make domain transfer and localized data assimilation possible, and fixes
the training and evaluation protocols used for every result in
Secs.~\ref{sec:results1d}--\ref{sec:failure}. Proofs of the
propositions are in Appendix~\ref{app:proofs}; numerical settings for
each system are in Appendix~\ref{app:architecture}.

%----------------------------------------------------------------------
\subsection{Setting and notation}
\label{sec:notation}
%----------------------------------------------------------------------

We consider a dissipative PDE $\partial_t u = N(u)$ for a scalar field
$u(\bm{x},t)$ on a periodic domain $\Omega$, either an interval of length
$L$ or a rectangle $L_x\times L_y$. The field is sampled on a uniform grid
of spacing $\Delta x$ and observed at snapshot interval $\Delta t$, giving
a discrete-time state $u_n = u(\cdot, n\Delta t)$. All fields are
standardized with a single mean and standard deviation computed on the
training data.

The domain is tiled by non-overlapping square \emph{patches} of
$p$ grid points per side, so that each patch has physical size
$h = p\,\Delta x$. The patch centres form a periodic \emph{site lattice}
$\Lat$ with $P = |\Lat|$ sites ($P = L/h$ in one dimension,
$P = (L_x/h)(L_y/h)$ in two). We write $\dist(j,k)$ for the periodic
(toroidal) distance between sites in units of the site spacing, and
$\Shift_s$ for the shift of a lattice field by $s\in\Z^{\dim}$ sites. A
shift of the physical field by $s$ whole patches is written
$\Tr_{s} u(\bm{x}) = u(\bm{x} - s h)$.

Two quantities set the scale of everything that follows. The first is
the number of active degrees of freedom, measured by the Kaplan--Yorke
dimension
\begin{equation}
  \DKY = k + \frac{\sum_{i=1}^{k}\lambda_i}{|\lambda_{k+1}|},
  \qquad
  k = \max\Big\{\, j : \textstyle\sum_{i=1}^{j}\lambda_i \ge 0 \Big\},
  \label{eq:dky}
\end{equation}
where $\lambda_1\ge\lambda_2\ge\cdots$ are the Lyapunov exponents
\citep{kaplanyorke1979}. The second is how $\DKY$ grows with domain
size. In \emph{extensive} chaos it grows in proportion to the volume,
\begin{equation}
  \DKY(\Omega) = \rho\,|\Omega| + o(|\Omega|),
  \label{eq:extensive}
\end{equation}
with a density $\rho$ that does not depend on $|\Omega|$
\citep{ruelle1982,manneville1985,egolfgreenside1994}. For the 1D
Kuramoto--Sivashinsky (KS) equation, $\rho\approx0.22$ per unit length
\citep{edson2019}. Equation~\eqref{eq:extensive} is the property a
surrogate must reproduce to be usable on domains larger than the one it
was trained on, and the property a fixed-size latent vector cannot
represent: its dimension is fixed at training time.

%----------------------------------------------------------------------
\subsection{Local latent fields}
\label{sec:local_field}
%----------------------------------------------------------------------

The latent state is a field on the site lattice,
\begin{equation}
  z = (z_j)_{j\in\Lat}, \qquad z_j \in \R^{c},
  \qquad \dlat = cP,
  \label{eq:latent_field}
\end{equation}
with $c$ channels per site. It is never flattened into an unstructured
vector except to call a library routine, and the flattening is always
site-major so that site identity is recoverable. The model has three
parts, each built only from operators that are applied identically at
every site and read a bounded neighbourhood.

\para{Encoder.}
The encoder $\E$ first maps each patch to a hidden vector by a strided
convolution whose kernel equals the patch (an exact, non-overlapping
patch embedding), then applies $n_e$ circular convolutions of radius
$r_e$ on the site lattice with GELU activations, and finally a
$1\times1$ convolution to $c-1$ channels. Channel $0$ is not learned: it
is the patch average of $u$,
\begin{equation}
  z_{j,0} = \frac{1}{|\text{patch } j|}\sum_{\bm{x}\in\text{patch } j} u(\bm{x}),
  \label{eq:anchor}
\end{equation}
which fixes one coordinate of the representation to a physical quantity
with the same meaning at every domain size. The output projection has
no bias: a per-channel offset shifts every site of the latent field by
the same constant, which the decoder can absorb and none of the training
losses constrains, so it drifts during training.

\para{Decoder.}
The decoder $\Dec$ mirrors the encoder: a $1\times1$ convolution from
$c$ channels to a hidden width, $n_d$ circular convolutions of radius
$r_d$, and a transposed convolution with kernel and stride equal to the
patch, which writes each site's hidden vector back onto its own patch.
The \emph{reach} of the encoder and decoder is $\Renc = r_e n_e$ and
$\Rdec = r_d n_d$ sites. When $\Renc=\Rdec=0$ the autoencoder is
\emph{per-patch}: site $j$'s latent is a function of patch $j$ alone,
and patch $j$ is reconstructed from $z_j$ alone. This is the canonical
autoencoder for the 2D systems (Sec.~\ref{sec:canonical_recipe}); the 1D
model uses $\Renc=\Rdec=6$ sites (Appendix~\ref{app:architecture}).

\para{Propagator.}
The latent dynamics are a residual map
\begin{equation}
  z^{n+1} = \Fp(z^{n}) = z^{n} + f_\theta(z^{n}),
  \label{eq:propagator}
\end{equation}
where $f_\theta$ (``local MLP'') embeds each site's $c$ channels to a
width $d_m$, applies $n_\ell$ blocks of the form
\begin{equation}
  \begin{aligned}
  a &\leftarrow a + \mathrm{Conv}_{w}\big(\mathrm{LN}(a)\big),\\
  a &\leftarrow a + \mathrm{MLP}\big(\mathrm{LN}(a)\big),
  \end{aligned}
  \label{eq:mixer_block}
\end{equation}
and projects back to $c$ channels per site. Here $\mathrm{Conv}_{w}$ is
a circular convolution of radius $w$ over the site lattice (kernel
$2w+1$ per axis), $\mathrm{LN}$ is layer normalization over channels at
each site, and $\mathrm{MLP}$ acts on each site independently. The token
mixer is an \emph{unconstrained linear} map of the neighbourhood: there
is no softmax normalization. Every self-attention propagator we tried
collapsed to a fixed point under free rollout (Sec.~\ref{sec:fail_attention});
replacing softmax attention by a linear mixer of the same radius removed
the collapse. The final projection is initialized to zero, so
$\Fp=\mathrm{id}$ at initialization and early rollouts are stationary
rather than divergent.

\para{Structural properties.}
All learned operators above are weight-shared and of finite radius on a
periodic lattice. Two consequences follow directly.

\begin{proposition}[Equivariance and size independence]
\label{prop:equivariance}
Let $\E$, $\Dec$ and $\Fp$ be built as above on a periodic lattice with at
least $2r+1$ sites along every axis on which a kernel of radius $r$ acts.
Then for every shift $s$ by a whole number of patches,
\begin{equation}
  \E\circ\Tr_{s} = \Shift_s\circ\E, \qquad
  \Dec\circ\Shift_s = \Tr_{s}\circ\Dec, \qquad
  \Fp\circ\Shift_s = \Shift_s\circ\Fp,
  \label{eq:equivariance}
\end{equation}
exactly, and the parameter vector $\theta$ of all three maps is the same
for every $P$.
\end{proposition}

\begin{corollary}[Transfer by adding sites]
\label{cor:transfer}
Hold $\Delta x$ and $p$ (hence $h$) fixed. Then a parameter vector
$\theta$ trained on any domain defines, without modification, a model on
every periodic domain whose side lengths are multiples of $h$, with
$P$ and $\dlat=cP$ growing in proportion to $|\Omega|$.
\end{corollary}

Corollary~\ref{cor:transfer} is the zero-shot transfer evaluated
throughout this paper: a checkpoint trained at one size is loaded at a
new $(N_x,P)$, with every weight copied unchanged. A global latent vector
admits no such map, because its bottleneck and propagator are tied to
one value of $\dlat$.

A local model is also extensive by construction, which explains both
what transfer can achieve and how it fails.

\begin{remark}[Local models are extensive]
\label{rem:extensive_model}
For translation-invariant lattice maps of finite range the Lyapunov
spectrum is expected, and found numerically, to have a thermodynamic
limit: the exponents, indexed per site, converge to a density as
$P\to\infty$
\citep{ruelle1982,livi1986,isola1990}; rigorous results exist for weakly
coupled maps \citep{bunimovichsinai1988}. For a local model this gives
$\DKY^{\text{model}}(\Omega) = \rho_\theta|\Omega| + o(|\Omega|)$ with a
density $\rho_\theta$ fixed by the trained weights, and the spectrum of a
small lattice already predicts that of a large one, as
\citet{carretero1999} showed for coupled map lattices. A local model can
therefore match the true density $\rho$. Its finite-size corrections,
however, are set by its own reach $\Rprop h$, not by the physics, so it
need not reproduce the true system's departures from
Eq.~\eqref{eq:extensive} on domains comparable to that reach. We use this
as an interpretive principle rather than a theorem about our networks: in
the results below the model's density is nearly constant across sizes,
and its largest errors occur where the domain is only a few reaches
wide.
\end{remark}

%----------------------------------------------------------------------
\subsection{Receptive fields and locality}
\label{sec:receptive}
%----------------------------------------------------------------------

The output of the encoder at site $j$ depends on patches within
$\Renc$ sites of $j$; one propagator step at site $j$ depends on latent
sites within $\Rprop = w\,n_\ell$ sites; the decoder output on patch $j$
depends on latent sites within $\Rdec$. After $k$ propagator steps,
the decoded field on patch $j$ depends on the initial field within
\begin{equation}
  R_k = \Renc + k\,\Rprop + \Rdec
  \label{eq:light_cone}
\end{equation}
sites. Unit tests confirm each reach on the implemented maps by masking
their Jacobians.

How large must $\Rprop$ be? For a dissipative PDE such as KS the exact
one-step map has an unbounded domain of dependence, but its influence
decays rapidly with distance, so the relevant quantity is an
\emph{interaction range}: the distance beyond which neighbouring values
no longer improve a one-step prediction. We measure it directly on the
1D latent. A single quadratic stencil, shared by all sites, that predicts
each site's one-step increment from the sites within a half-width $w$
explains $58\%$ of the variance at $w=1$, $82\%$ at $w=2$, $88\%$ at
$w=3$ and saturates at $90\%$ from $w=4$ on, the value it also reaches
with access to the whole ring (a linear stencil explains at most
$2.5\%$ at any width; Fig.~\ref{fig:appendix_panels}(c)). The interaction range of the 1D dynamics at
$\Delta t=1$ is therefore three to four sites, $19$--$25$ length units, or
two to three pattern wavelengths. The trained propagator's reach of six
sites ($37.5$ units) covers it with a margin; the 2D models' six-site
reach spans about four pattern wavelengths in 2D KS ($33$ units) and
three forcing wavelengths in Kolmogorov flow ($4.7$ units).

A second, more familiar length is set by the speed at which
perturbations spread. In spatiotemporal chaos this speed is finite and
is characterized by the largest velocity $v_*$ with a positive comoving
(velocity-dependent) Lyapunov exponent
\citep{kaneko1986,deisslerkaneko1987}. For 1D KS we measure
$v_*=1.26\pm0.04$ length units per time unit (150 trajectories), so a
perturbation spreads $v_*\Delta t\approx1.3$ units, a fifth of a site,
during one snapshot interval. This is a lower bound on the reach a
propagator needs, and it is far from binding here: the interaction range
is set by how far the one-step map couples values, not by how far
perturbations travel. The same distinction applies to localization in
data assimilation (Sec.~\ref{sec:da_method}): an analysis increment at
site $j$ must reach every site that the decoder or the forecast couples to
$j$. With a per-patch decoder the decoder contributes nothing, and the
light-cone term $v_*\Delta t_{\mathrm{win}}$ over an assimilation window
of five snapshots is about $6$ length units in 1D, far inside the
localization radii we use.

%----------------------------------------------------------------------
\subsection{Training on windows of a large domain}
\label{sec:crop_training}
%----------------------------------------------------------------------

A model trained on a small periodic box sees only that box's statistics,
and a propagator whose reach covers most of the box can exploit the
periodic wrap-around, a self-interaction that does not exist on a larger
domain. In our 2D experiments this happened: a propagator with
$\Rprop=6$ trained on an $8$-site periodic ring covers $75\%$ of it, and
such models fail to transfer (Sec.~\ref{sec:box_vs_window}). We instead
simulate one large periodic domain and train only on random rectangular
windows (\emph{crops}) of it. A crop has no boundary condition: its
edges are open, and the training losses are restricted to the sites a
light-cone argument certifies.

\para{Stage A: autoencoder.}
A crop of $W$ sites per cropped direction is encoded and decoded. Sites
within $\Renc+\Rdec$ of a crop edge are excluded from the
reconstruction loss, because padding at the edge reaches them. For a
per-patch autoencoder this margin is zero.

\para{Stage B: propagator with a forced halo.}
The trained encoder is frozen and applied to many crops of every training
trajectory at every time, giving a bank of latent crop trajectories.
The outer $\Renc$ sites per cropped side are dropped, since their
encoding saw padding. The remaining strip is split into a \emph{halo}
of $\eta$ sites per side and a \emph{free interior} $I$. A training
rollout starts from the true encoding, applies the propagator, and after
every step overwrites the halo with the true encoding at that time; only
the interior is scored. This is the construction of a limited-area
forecast model driven by lateral boundary data from a parent run.

\begin{proposition}[Forced-halo exactness]
\label{prop:halo}
Let $\eta\ge\Rprop$. Then for every rollout length $k$, the forced-halo
prediction on the free interior equals the prediction of the same
weights on any larger periodic domain containing the crop, in which every
site outside $I$ is set to its true value at every step. If
$\eta<\Rprop$, the equality fails in general.
\end{proposition}

Proposition~\ref{prop:halo} says the crop loss is the loss of a
well-defined object, the model restricted to $I$ with exact lateral
boundary data, and that no padding artifact enters the gradient. It is
checked numerically (agreement to $10^{-6}$ for strips cropped in one
direction and $10^{-5}$ for windows cropped in both, and companion tests
that halos of $\Rprop-1$ and $\Rprop-2$ break it). The one
mismatch between training and deployment is that interior sites adjacent
to the halo see true neighbours in training and predicted neighbours on
the periodic deployment domain; it is confined to a strip of width
$\Rprop$.

For 2D Kuramoto--Sivashinsky, which is isotropic, crops are cut in both
directions, so no training sample contains a whole periodic direction.
For Kolmogorov flow, whose forcing depends on $y$, crops are strips in
$x$ that keep the full periodic height, and transfer grows the domain in
$x$ only.

%----------------------------------------------------------------------
\subsection{Training objective}
\label{sec:objective}
%----------------------------------------------------------------------

\para{Stage A.}
For a batch of crops with interior latents $\{z_j\}$ pooled over sites
and batch members, let $C = \Cov(z)\in\R^{c\times c}$ be their channel
covariance. The loss is
\begin{equation}
\begin{split}
  \mathcal{L}_A ={}&
  \big\|\,\Dec(\E(u)) - u\,\big\|^2_{I}
  + \wvar\,\frac1c\sum_{k=1}^{c}\big(C_{kk}-1\big)^2\\
  &- \wld\,\frac1c\log\det\!\big(C + \epsilon I_c\big),
\end{split}
  \label{eq:stageA}
\end{equation}
where $\|\cdot\|_I$ is the mean squared error over the unmasked interior.
The two regularizers do different jobs. The log-determinant barrier
penalizes any small eigenvalue of $C$, whether it comes from a channel
with small variance or from a channel that is nearly a linear combination
of others, so it prevents correlated (rank) collapse that a per-channel
penalty cannot see. It does not, however, constrain the overall scale:
it decreases when all eigenvalues grow together, and the encoder and
decoder can rescale against each other with no change in reconstruction.
The variance term anchors that scale. With $\wvar=0$ we observed the
variance diagnostic grow from $8\times10^{2}$ to $2.5\times10^{8}$ over
training while the reconstruction error stayed flat.

\begin{remark}[The barrier offset $\epsilon$ is not a numerical detail]
\label{rem:epsilon}
The barrier's force on an eigenvalue $\mu$ of $C$ is
$-\partial_\mu\log(\mu+\epsilon) = -1/(\mu+\epsilon)$. At initialization
the smallest eigenvalue of $C$ for our architecture is
$0.8$--$4.4\times10^{-3}$ (ten random seeds), so changing $\epsilon$
from $10^{-3}$ to $0$ changes the force on the weakest directions by a
factor of order one. Section~\ref{sec:fail_mse} shows that this choice
changes the trained model's dimension at large domain sizes as much as
the rollout length does.
\end{remark}

\para{Stage B.}
The propagator is trained on forced-halo rollouts of length $k$,
\begin{equation}
  \mathcal{L}_B = \frac1k\sum_{m=1}^{k}
  \big\|\hat z^{m} - z^{m}\big\|_{I}^{2},
  \quad \hat z^{0} = z^{0} + \xi,\ \xi\sim\mathcal N(0,\sigma^2 I),
  \label{eq:stageB}
\end{equation}
where $\hat z^{m}$ is the forced-halo prediction. The rollout length $k$
is ramped from $2$ to an intermediate value $\kmid$ over the first $58\%$
of training, then to its final value $\kmax$, and held there for the last
$30\%$; the input noise $\sigma$ is
annealed linearly from $0.10$ to $0.02$. Long free rollouts are needed:
short ones leave the free-running model under-constrained and produce
spurious, initial-condition-dependent regimes (Sec.~\ref{sec:fail_regimes}).
But long rollouts scored by mean squared error carry a known bias.

\begin{lemma}[Squared-error rollouts favour contraction]
\label{lem:mse}
Let $(z^0, z^k)$ be jointly distributed with finite second moments, and
let $\tilde z^0$ be the (possibly noisy) model input. Among all
measurable maps $\Phi$, the expected loss
$\Exp\|\Phi(\tilde z^0)-z^k\|^2$ is minimized by
$\Phi^\star(\tilde z^0)=\Exp[z^k\mid\tilde z^0]$, and
\begin{equation}
  \tr\Cov\big(\Phi^\star(\tilde z^0)\big)
  = \tr\Cov(z^k) - \Exp\,\tr\Cov\big(z^k\mid\tilde z^0\big).
  \label{eq:total_variance}
\end{equation}
\end{lemma}

The second term in Eq.~\eqref{eq:total_variance} is the forecast
uncertainty at lead $k$. In a chaotic system it is non-zero at every
$k>0$: input noise of variance $\sigma^2$ grows along the leading
direction as $\sigma^2 e^{2\lambda_1 k\Delta t}$ until it saturates at the
climatological variance. A propagator trained toward $\Phi^\star$ is
therefore rewarded for rollouts whose spread is smaller than that of
real trajectories, that is, for contracting along the most unstable
directions, and the reward grows with $k$. Our training horizons are
short in Lyapunov units: $\kmax\Delta t\,\lambda_1\approx1.1$ for 1D KS,
$0.5$--$0.65$ for 2D KS, and $0.2$--$1.0$ for Kolmogorov flow across the
three drags. The effect is therefore a bias that grows with $\kmax$,
not the collapse to the climatological mean that a horizon of many
Lyapunov times would produce. It is the mechanism behind the
``blurring'' of machine-learned weather models
\citep{lam2023graphcast,bonavita2024}, and for a learned surrogate it
predicts an underestimate of $\DKY$ that grows with $\kmax$, which is
what we observe (Sec.~\ref{sec:fail_mse}).

\begin{remark}[Input noise penalizes growth directly]
\label{rem:noise_penalty}
Training with input noise is, to leading order, a penalty on the
Jacobian of the trained map \citep{bishop1995}. Expanding
$\Phi(z+\xi)$ for small $\xi\sim\mathcal N(0,\sigma^2 I)$,
\begin{equation}
\begin{split}
  \Exp_\xi\big\|\Phi(z+\xi)-y\big\|^2
  ={}& \big\|\Phi(z)-y\big\|^2 + \sigma^2\big\|\partial_z\Phi(z)\big\|_F^2\\
  &+ \sigma^2\,\big(\Phi(z)-y\big)^{\!\top}\Delta_z\Phi(z) + O(\sigma^4),
\end{split}
  \label{eq:noise_penalty}
\end{equation}
where $\Delta_z\Phi$ is the vector of Laplacians of the components of
$\Phi$; the last term averages to zero at the regression optimum. For the
$k$-step map the Frobenius norm is at least the largest finite-time
growth factor, $\|\partial_z\Phi\|_F^2\ge e^{2\lambda^{(k)}_1 k\Delta t}$, so
input noise penalizes exactly the expansion that produces positive
exponents, with a strength $\sigma^2$. Rollout length and input noise
are two handles on one bias. The one run that raised both (rollouts to
$\kmax=32$ with terminal noise $0.05$ instead of $0.02$,
Sec.~\ref{sec:fail_regimes}) underestimates $\DKY$ by $4\%$, $13\%$ and
$14\%$ at $1\times$, $2\times$ and $3\times$.
\end{remark}

Rollouts that are too short fail in the opposite way
(Sec.~\ref{sec:fail_regimes}). We treat $\kmax$ and the noise schedule as
a trade-off between these two failure modes, not as quantities to
maximize.

%----------------------------------------------------------------------
\subsection{The canonical recipe}
\label{sec:canonical_recipe}
%----------------------------------------------------------------------

All 2D results in the main text share one architecture and training
procedure: a per-patch autoencoder ($\Renc=\Rdec=0$) with $c=6$ channels
(five learned and the anchor) and hidden width $32$; a local-MLP
propagator with $n_\ell=2$ blocks of radius $w=3$ ($\Rprop=6$ sites),
token width $d_m=32$ and MLP ratio $4$; training on crops of one large
domain with a forced halo $\eta=6$; and Stage~A for $8000$ steps and
Stage~B for $6000$--$9000$ steps with AdamW, linear warm-up and cosine
decay. The models differ in the rollout length $\kmax$ and in the
regularizer weights. $\kmax=32$ (2D KS and Kolmogorov flow at
$\alpha=0.025$) is the long-rollout setting introduced to suppress
spurious free-running regimes (Sec.~\ref{sec:fail_regimes}); the
$\alpha=0.05$ model is the earlier, independently verified $\kmax=12$
model, which a $\kmax=32$ retraining does not improve
(Sec.~\ref{sec:canonical_results}); and at $\alpha=0.1$, where $\kmax$ and
the barrier offset $\epsilon$ interact (Sec.~\ref{sec:fail_mse}), the
model ($\kmax=16$) was chosen among six variants by a rule that uses no
ground truth (Sec.~\ref{sec:canonical_results}).

The 1D model predates this recipe. It was trained on whole periodic
domains of $16$ sites with a full-reach encoder and decoder
($\Renc=\Rdec=6$, $c=3$), jointly with a short-horizon auxiliary
propagator, and with an additional penalty that favours spatial
coherence of the latent across neighbouring sites; the propagator was
then refined with rollouts up to $\kmax=12$. It transfers without crop
training (Sec.~\ref{sec:results1d}), so we did not retrain it with the 2D
recipe. Every setting of every model is listed in
Appendix~\ref{app:architecture}.

%----------------------------------------------------------------------
\subsection{Evaluation}
\label{sec:evaluation}
%----------------------------------------------------------------------

\para{Lyapunov spectra.}
We compute Lyapunov spectra by the Benettin algorithm with QR
reorthonormalization \citep{benettin1980}. For the true PDE, tangent
vectors are propagated with the exact Fr\'echet derivative of the
time-stepper in double precision. For a model, tangent vectors are
propagated through $\Fp$ by forward-mode automatic differentiation in
single precision and re-orthonormalized every 10 steps by a
double-precision QR on the CPU. The exponents are reported per unit
physical time, $\lambda_i = (n\Delta t)^{-1}\sum_{q}\log|R^{(q)}_{ii}|$,
where $R^{(q)}$ is the triangular factor of the $q$-th reorthonormalization
and $n$ the number of model steps. From each spectrum we report
$\lambda_1$, the number of positive exponents $n_+$, and $\DKY$. Each
model spectrum uses all $\dlat$
tangent directions, a warm-up of 200 steps and $n=2000$ steps (2000 time
units in 1D, 200 in 2D KS, 400 in Kolmogorov flow); the computation
raises an error if the cumulative sum in Eq.~\eqref{eq:dky} never turns
negative, so $\DKY$ is always bracketed rather than extrapolated. The
spectrum is that of the latent map; comparing it with the PDE's assumes
that the decoder is a smooth, injective map of the attractor, which the
agreement of the spectra themselves supports (Sec.~\ref{sec:results1d}).

A model trajectory starts from the encoding of a true on-attractor
state, standardized with the training statistics, and each model is
evaluated from three independent initial conditions per domain size
unless stated otherwise. A finite-time estimate from one initial
condition is a random quantity, for the model as for the truth: two runs
that differ only at round-off level (single versus double precision)
follow different chaotic trajectories. Averaged over three initial
conditions, double-precision propagation changes $\DKY$ by at most
$0.3\%$ in 1D and 2D KS and by $2.2\%$ in Kolmogorov flow at the native
size, where single runs scatter by $25\%$ across initial conditions; we
therefore report means over initial conditions and treat single-run
differences below that scatter as noise.

\para{Ground truth and its uncertainty.}
True spectra are computed after a spin-up of $100$--$200$ time units and
a warm-up of the tangent basis. For Kolmogorov flow we average over $120$
time units and record the running estimate of $\DKY$ at intermediate
times to check convergence; the window used for every truth value is
listed in Appendix~\ref{app:groundtruth}. The truth's finite-time $\DKY$ is
random in the same way as the model's. In Kolmogorov flow its
seed-to-seed spread is $4$--$13\%$ at the sizes we use
(Sec.~\ref{sec:truth_uncertainty}), comparable to many model errors, so we
report 2D results as scaling with domain size, with truth error bars,
rather than as single-point matches.

\para{Free-rollout statistics.}
Before any Lyapunov computation, each model is run freely for several
hundred time units and its field variance and energy spectrum, split
into large-scale and forcing-scale bands, are compared with the training
data or with the truth.
This check is cheap and detects drifting or multistable models that a
short teacher-forced validation loss does not (Sec.~\ref{sec:fail_regimes}).

%----------------------------------------------------------------------
\subsection{Latent ensemble data assimilation}
\label{sec:da_method}
%----------------------------------------------------------------------

Observations are taken of the physical field,
$y_n = \Hobs(u_n) + \nu_n$, $\nu_n\sim\mathcal{N}(0,R)$, where $\Hobs$
selects grid points on a regular sparse network. The filter state is
an ensemble $\{z^{(e)}\}_{e=1}^{\Nens}$ of latent fields, forecast with
$\Fp$ plus additive model noise, and the observation operator in latent
space is $h = \Hobs\circ\Dec$. We use three analysis schemes:
\begin{itemize}
\item a particle flow filter (the Gauss--Newton ``NAT'' variant of
  \citet{huvanleeuwen2021}), used in 1D;
\item a stochastic ensemble Kalman filter with perturbed observations
  \citep{burgers1998,houtekamermitchell1998} and Schur-product
  covariance localization \citep{houtekamermitchell2001}, with gain
  $K = (\rho\circ B)J^{\top}\big(J(\rho\circ B)J^{\top}+R\big)^{-1}$,
  where $B$ is the ensemble covariance, $\rho$ the taper, and $J$ the
  Jacobian of $h$ at the ensemble mean; used in 1D and 2D;
\item a per-site local ensemble Kalman filter in the local-analysis
  form of \citet{hunt2007}, with observation-space (R-)localization and
  the same perturbed-observation update: each site is updated from the
  observations within the taper's support, with each observation's error
  variance divided by its taper value; used in 2D.
\end{itemize}
The taper is the Gaspari--Cohn function \citep{gaspari1999} of the site
distance $\dist(j,k)$, with toroidal distance in 2D. As a distance-free
baseline we use sampling error correction (SEC), which replaces the
taper by a correction factor fitted offline from the sample correlation
\citep{anderson2012}.

The per-patch decoder makes the observation operator exactly local in
latent space.

\begin{proposition}[Block-diagonal observation Jacobian]
\label{prop:blockdiag}
Let $\Rdec=0$ and let $\Hobs$ select grid points. Then
$\partial h_i/\partial z_{j,k} = 0$ for every channel $k$ unless
observation $i$ lies in patch $j$. Consequently:
\begin{enumerate}
\item[(a)] the full Jacobian $J\in\R^{m\times\dlat}$ is recovered exactly
  from $c$ forward-mode Jacobian--vector products with the tangents
  $v^{(k)} = \sum_{j\in\Lat} e_{j,k}$, $k=1,\dots,c$, for any $P$ and any
  number of observations $m$;
\item[(b)] in the per-site filter with R-localization, with either the
  perturbed-observation or the ensemble-transform update, the analysis at
  site $s$ depends only on the ensemble at the sites that own an
  observation within the taper support of $s$, and no row of $J$ is
  truncated in computing it.
\end{enumerate}
\end{proposition}

Part (a) is a column colouring of $J$ in the sense of
\citet{curtis1974}: all columns belonging to the same channel can share
one colour, because no row has two non-zeros in such a group. With
$\Rdec>0$, row $i$ would have non-zeros at every site within $\Rdec$ of
its patch, the colouring would need $c(2\Rdec+1)^{\dim}$ colours, and the
per-site analysis in (b) would have to include those sites as well or
truncate $h$. Neither the cost of the colouring nor the memory of the
per-site filter depends on the domain size; the dense covariance-localized
filter, which forms $\dlat\times\dlat$ matrices, scales as $P^2$ in memory.

\para{How large an ensemble is needed.}
An ensemble filter needs enough members to span the directions along
which forecast errors grow. Without localization these are the unstable
directions of the whole attractor, so $\Nens$ must grow with $\DKY$ and
therefore with $|\Omega|$. With localization, each local analysis only
needs to span the unstable directions active inside the taper support
$A_{\mathrm{loc}}$ \citep{patil2001,ott2004}, which suggests
\begin{equation}
  \Nens \gtrsim \rho\,|A_{\mathrm{loc}}|,
  \label{eq:nens_local}
\end{equation}
independent of the total domain size. For 1D KS with the taper used in
Sec.~\ref{sec:da_results} (support radius $6$ sites of length $6.25$),
Eq.~\eqref{eq:nens_local} gives $\Nens\gtrsim 0.22\times75\approx17$.
Combined with Corollary~\ref{cor:transfer}, this is the reason to expect a
transferred local model to keep a fixed-size filter useful as the domain
grows, while an unlocalized filter of the same size degrades.
